% ============ 参考文献 (原书 pp.156) + 作者简介 (原书卷末) ============

\chapter*{参考文献}

就日语教材而言, [2], [3] 是新刊. 这两本书是长达 600 页的巨著, 在读完本书内容之后再读为好. [4]-[6] 也可以说是与 [2], [3] 不同的, 从本书延伸的文献. 就英语教材而言, 有 [7]-[9], [12] 等. [7] 是由约 200 页构成的相对较短的书籍, 但在读同样由两部构成的 [8], [9] 以及巨著 [12] 之前, 先读这本书的话, 虽然无法概观这些巨著, 却能体会其工作的大致氛围. 数列空间中的调和分析请参考 [1]. 另外, 练习问题 3.5 与 5.10 分别采录自 [10], [11].

\begin{itemize}
\item[{[1]}] 新井仁之, フーリエ解析学 (講座 数学の考え方 17), 朝倉書店, 2003.
\item[{[2]}] 宮地晶彦, ユークリッド空間上のフーリエ解析 I (朝倉数学大系 13), 朝倉書店, 2021.
\item[{[3]}] 宮地晶彦, ユークリッド空間上のフーリエ解析 II (朝倉数学大系 14), 朝倉書店, 2021.
\item[{[4]}] 藪田公三, 佐藤圓治, 田中仁, 宮地晶彦, 中路貴彦, 古典調和解析 (解析学百科 I), 朝倉書店, 2008.
\item[{[5]}] 藪田公三, 特異積分 (岩波数学叢書), 岩波書店, 2010.
\item[{[6]}] 澤野嘉宏, ベゾフ空間論, 日本評論社, 2011.
\item[{[7]}] J. Duoandikoetxea, \emph{Fourier Analysis}, translated and revised from the 1995 Spanish original by D. Cruz-Uribe, Graduate Studies in Mathematics, 29, American Mathematical Society, Providence, RI, 2001.
\item[{[8]}] L. Grafakos, \emph{Classical Fourier analysis}, Graduate Texts in Mathematics, 249, Springer, New York, 2014.
\item[{[9]}] L. Grafakos, \emph{Modern Fourier analysis}, Graduate Texts in Mathematics, 250, Springer, New York, 2014.
\item[{[10]}] V. Burenkov, A. Ghorbanalizadeh and Y. Sawano, Uniform boundedness of Kantorovich operators in Morrey spaces, \emph{Positivity}, 22, 1097--1107 (2018).
\item[{[11]}] K. Morii, T. Sato and Y. Sawano, Certain identities on derivatives of radial homogeneous and logarithmic functions, \emph{Commun.\ Math.\ Anal.}, 9, no.\ 2, 51--66 (2010).
\item[{[12]}] E.M. Stein, \emph{Harmonic Analysis, Real-variable Methods, Orthogonality, and Oscillatory Integrals}, Princeton University Press, Princeton, NJ, 1993.
\end{itemize}

\clearpage
\chapter*{作者简介}

\begin{center}
\large 泽野嘉宏 \quad (さわのよしひろ)
\end{center}

\begin{itemize}
\item 1979 年 生于静冈县
\item 2006 年 东京大学大学院数理科学研究科博士后期课程修了, 博士(数理科学)
\item 2008 年 学习院大学理学部助教
\item 2009 年 京都大学理学部助教
\item 2012 年 首都大学东京理学部准教授
\item 2020 年 中央大学理工学部教授, 至今
\end{itemize}

2009 年获日本数学会奖建部贤弘特别奖.

专业, 研究方向: 调和分析, 函数空间论.

主要著作: \emph{Theory of Besov spaces}, (Springer, 2018); \emph{Morrey Spaces: Introduction and Applications to Integral Operators and PDE's}, Volume I, II, (Chapman and Hall/CRC, 2020); 『ベゾフ空間論』(日本评论社, 2011), 等.
